% Nouveau EGA V (cote 25): archive cote 25, site page 145 (PDF page 146).
% English translation of the checked French annotated master.
% AI-translated by GPT-6 Astra (OpenAI, ChatGPT Pro mode).
\scanpage{25/145}
\mspage{\hand{6}}

\uline{Proposition} 7. Under the preceding conditions (to be recalled), the underlying
set of $V(\hand{\phi})^{\text{sup sing}}$ is precisely the set of points of
$V(\hand{\phi})^{\text{sing}}$ \var{\xout{\unsure{mixed}}}{}{typewriter strike-out}{ega25_25_145-v0001} where the morphism $V(\hand{\phi})^{\text{sing}}\hand{\longrightarrow} Y$ (of preschemes smooth
over S of relative dimensions m-1 and m respectively) is \uline{ramified}.

In papa's terms (which should be pointed out in a remark), a point
$x \hand{\in} V(\hand{\phi})$ is therefore supersingular relative to Y iff ``it consists of at
least two coincident singular points'' ....

One can and should make Prop. 7 more precise from the point of view \var{\inserted[typed]{of an identity of}}{of an identity of}{insertion}{ega25_25_145-v0002} \var{\xout{of the}}{}{typewriter strike-out}{ega25_25_145-v0003} subpreschemes,
and not only of subsets. Indeed, $V(\hand{\phi})^{\text{sup sing}}$ was
defined as a closed subprescheme of X, and one can also define
a natural closed subprescheme of $V^1$ in such a way that its underlying
set is the set of ramification points relative to Y.
Indeed, one must express the set of points where a certain homomorphism
of loc free modules $Q$ ( $=\underline{\hand{\Omega}}^1_{Y/S}\otimes\underline{O}_{V^1}$ ) $\hand{\longrightarrow}$ $R$ ( $=\underline{\hand{\Omega}}^1_{V^1/S}$ )
is not surjective. If q and m are their respective ranks, this is also
the set of points where $\hand{\bigwedge}^{m} Q \hand{\rightarrow} \hand{\bigwedge}^{m} R$ is not surjective, which is
also the set of zeros of the evident section of \uline{Hom} ( $\hand{\bigwedge}^{m} Q$, $\hand{\bigwedge}^{m} M$)
$\hand{\simeq} (\hand{\bigwedge}^{q-r} Q)\otimes(\hand{\bigwedge}^{r} R^{\hand{\illegible{1}}})$$\otimes(\hand{\bigwedge}^{q} Q)$\hand{$^{\vee}$}, hence the underlying set of the closed \var{\inserted[typed]{sub}}{sub}{insertion}{ega25_25_145-v0004}prescheme
of zeros of this section, namely \uline{Ram}($V^1/Y$). I claim that this latter
subscheme is identical to $V(\hand{\phi})^{\text{sup sing}}$. This is a simple exercise
on the diagram above, taking into account that $V(\hand{\phi})^{\text{sup sing}}$ is defined
by the same procedure as that given for $Q\hand{\longrightarrow} R$, but in terms of
the homomorphism $P$ ( $= \overset{\hand{\vee}}{\underline{P}}{}^1_{X/Y}\otimes\underline{O}_{V^1}$ ) $\hand{\longrightarrow}$ $S$ ($=\underline{\hand{\Omega}}^1_{X/Y}\otimes\underline{O}_{V^1}$), as
follows from the description of \hand{$\mu$} given above. We are then
reduced to the following general situation:

On a \var{\deleted{prescheme}}{}{deletion}{ega25_25_145-v0006} \var{\inserted{ringed space}}{ringed space}{insertion}{ega25_25_145-v0007} W we have a locally free module \var{\inserted[typed]{M}}{M}{insertion}{ega25_25_145-v0008} of rank m,
and two locally free submodules P and Q of respective ranks

\mspage{\hand{145}}
\tnote{25-145-1}{Papa's terms renders termes de papa literally, preserving the author's informal phrase. Loc abbreviates locally. The source ramffié is translated as ramified. The n overwritten by n in On has no English effect. In respectifs the overwritten underletter is doubtful p/s, while the final p is secure; the English adjective is certain, and the source-letter uncertainty is recorded here. The source alternates q and m with q and r in the rank and exterior-power expressions and has M in the second Hom argument; all these mathematical tokens are retained.}
